\section{Limitations and Conclusion}
\label{sec:conclusion}

\paragraph{Limitations.}
Our fidelity evaluation relies mainly on synthetic needles inserted into natural text; LongBench establishes that the method does no harm on natural documents but not how much it helps on their identifiers, for lack of a benchmark that asks for verbatim copies of naturally occurring strings.
We evaluate black-box readers only and therefore cannot say how enlarged glyphs change a model's internal attention.
Budgets follow one patch-based accounting rule (\cref{eq:vtok}); providers tile images differently, and while billed costs confirm the savings for the readers we tested, other tokenizers may shift them.
The detector is rule-based: it is fail-safe but recall-bounded, and hazards it does not recognize receive no protection.
Finally, reallocation alone raises recall without reducing the share of confident near-misses; only anchors remove them.

\paragraph{Conclusion.}
Optical context compression decides how many pixels a document receives but not where they go, and its failures concentrate on exactly the strings that must survive verbatim.
Because the renderer reads the text before it draws it, it can spend resolution where errors hurt, reinvest the savings in anchors, and fall back to text when a document is too dense with identifiers.
These steps are training-free and model-agnostic, and across six VLMs they make optical compression reliable for embedded identifiers: on long documents with few identifiers, the routed system retrieves them nearly as well as plain text at about one fifth of its cost, and it sends identifier-dense documents to text.
